% ============================================================
% 第7章补充：经典例题与自学元素
% 插入位置：第7章现有miniquiz之后（\end{miniquiz} 之后）
% ============================================================

\section*{习题精选与详解}
\addcontentsline{toc}{section}{习题精选}

\begin{example}{麦克斯韦分布验证Vlasov方程的稳态解}{maxwell-vlasov-verify}
设一维平衡分布为麦克斯韦分布：
\[
f_0(v) = n\sqrt{\frac{m}{2\pi k_{\!B}T}}\,\exp\!\left(-\frac{mv^2}{2k_{\!B}T}\right)
\]
其中 $n$、$T$ 为常数，电场 $\vect{E}=\vect{0}$，磁场 $\vect{B}=\vect{0}$。验证 $f_0(v)$ 是Vlasov方程的解，即分布函数不随时间演化。
\tcblower
\textbf{解答：}

一维Vlasov方程（无电磁场）为：
\[
\frac{\partial f}{\partial t} + v\frac{\partial f}{\partial x} + \frac{qE}{m}\frac{\partial f}{\partial v} = 0
\]

将 $f = f_0(v)$ 代入，逐项计算：

\textbf{第一项}：$\frac{\partial f_0}{\partial t}$。由于 $f_0$ 只依赖于速度 $v$，不显含时间 $t$，故
\[
\frac{\partial f_0}{\partial t} = 0
\]

\textbf{第二项}：$v\frac{\partial f_0}{\partial x}$。由于 $f_0$ 只依赖于速度 $v$，不显含空间坐标 $x$，故
\[
\frac{\partial f_0}{\partial x} = 0 \quad\Rightarrow\quad v\frac{\partial f_0}{\partial x} = 0
\]

\textbf{第三项}：$\frac{qE}{m}\frac{\partial f_0}{\partial v}$。题设 $E=0$，故此项为零。

三项之和为零：
\[
0 + 0 + 0 = 0
\]

因此 $f_0(v)$ 满足Vlasov方程。\textbf{物理意义}：在没有电磁场和梯度的情况下，任何只依赖 $v$ 的分布都是Vlasov方程的稳态解——麦克斯韦分布只是其中一员。

\begin{warning}{"稳态解"不等于"热平衡"，也不等于"稳定"}{vlasov-steady-caution}
这是本节最容易滑过去的一步。上面的验证只说明 $f_0(v)$ 满足 $\partial_tf=0$，它\emph{没有}说明系统通过无碰撞过程达到了热平衡。请把三件事分清：
\begin{itemize}
\item \textbf{稳态}：$\partial f_0/\partial t=0$，分布不随时间变。这只需要 $f_0$ 不含 $t$、$x$，\emph{任何} $f_0(v)$ 都满足。
\item \textbf{热平衡}：需要由\emph{碰撞}（或某种耗散）驱动到麦克斯韦分布。纯 Vlasov 方程（碰撞项为零）里根本没有把分布"拉向"麦克斯韦的机制，所以"代入麦克斯韦分布能验证稳态"\emph{不能}证明无碰撞热化。
\item \textbf{稳定}：还需要分布满足无自由能的判据。一个稳态分布完全可以是\emph{不稳定}的——见下面的反问题。
\end{itemize}

\textbf{反问题（请自行验证）：} 取双峰分布 $f_0(v)=n[A_1\ee^{-(v-u)^2/2v_{t1}^2}+A_2\ee^{-(v+u)^2/2v_{t2}^2}]$（两束相向运动的电子），它同样只依赖 $v$，因此\emph{同样}是 Vlasov 方程的稳态解——但它是\textbf{双流不稳定性}的经典不稳定平衡（第9章）。可见"满足 Vlasov 稳态"与"能稳定存在"是两回事。真正决定稳定性的是 $\partial f_0/\partial v$ 的符号结构（第9章、本章末的 bump-on-tail），而不是"是否满足 $\partial_tf=0$"。
\end{warning}
\end{example}

\begin{example}{Landau阻尼率的数值估算}{landau-damping-calc}
实验室氩等离子体中，测得电子密度 $n_e = 10^{18}\,\mathrm{m}^{-3}$，电子温度 $T_e = 4\,\mathrm{eV}$。一列电子Langmuir波的波数为 $k = 10^{4}\,\mathrm{m}^{-1}$。请估算该波的Landau阻尼率 $|\omega_i|$ 和对应的衰减时间 $\tau = 1/|\omega_i|$。

\textbf{公式的适用条件（先读）：} 本题所用的阻尼率公式
\[
\omega_i = -\sqrt{\frac{\pi}{8}}\,\frac{\omega_{pe}}{(k\lambda_D)^3}\exp\!\left(-\frac{1}{2k^2\lambda_D^2} - \frac{3}{2}\right)
\]
是\textbf{长波、弱阻尼的渐近式}，成立条件是 $k\lambda_D\ll1$ 且 $|\omega_i|\ll\omega_r$。它由 $\omega$ 的虚部在实轴附近做展开得到，一旦 $k\lambda_D$ 接近或超过 1，该展开失效，必须改解含等离子体色散函数 $Z(\zeta)$ 的复根。下面先算一个\textbf{确实落在适用区}的例子，再专门讨论越界会怎样。
\tcblower
\textbf{解答：}

\textbf{步骤1：计算德拜长度。}
已知 $k_{\!B}T_e = 4\,\mathrm{eV} = 6.408\times 10^{-19}\,\mathrm{J}$，
\[
\lambda_D = \sqrt{\frac{\eps_0 k_{\!B}T_e}{n_e e^2}}
= \sqrt{\frac{8.854\times 10^{-12}\times 6.408\times 10^{-19}}{10^{18}\times (1.602\times 10^{-19})^2}}
\approx 1.49\times 10^{-5}\,\mathrm{m} = 14.9\,\mu\mathrm{m}
\]

\textbf{步骤2：计算无量纲参数 $k\lambda_D$。}
\[
k\lambda_D = 10^{4}\times 1.49\times 10^{-5} \approx 0.149 \ll 1 \quad\checkmark
\]
落在长波渐近式适用区。

\textbf{步骤3：代入Landau阻尼率公式。}
由教材公式：
\[
\omega_i = -\sqrt{\frac{\pi}{8}}\,\frac{\omega_{pe}}{(k\lambda_D)^3}\,
\exp\!\left(-\frac{1}{2k^2\lambda_D^2} - \frac{3}{2}\right)
\]

先计算电子等离子体频率：
\[
\omega_{pe} = \sqrt{\frac{n_e e^2}{\eps_0 m_e}}
= \sqrt{\frac{10^{18}\times (1.602\times 10^{-19})^2}{8.854\times 10^{-12}\times 9.109\times 10^{-31}}}
\approx 5.64\times 10^{10}\,\mathrm{rad/s}
\]

计算指数部分：
\[
\frac{1}{2k^2\lambda_D^2} = \frac{1}{2\times (0.149)^2} = \frac{1}{0.0444} \approx 22.5
\]
\[
-\frac{1}{2k^2\lambda_D^2} - \frac{3}{2} = -22.5 - 1.5 = -24.0
\]
\[
\exp(-24.0) \approx 3.78\times 10^{-11}
\]

计算前置因子：
\[
\sqrt{\frac{\pi}{8}}\,\frac{\omega_{pe}}{(k\lambda_D)^3}
= 0.627\times \frac{5.64\times 10^{10}}{(0.149)^3}
= 0.627\times \frac{5.64\times 10^{10}}{3.31\times 10^{-3}}
\approx 1.07\times 10^{13}\,\mathrm{rad/s}
\]

\textbf{步骤4：得到阻尼率。}
\[
|\omega_i| = 1.07\times 10^{13}\times 3.78\times 10^{-11} \approx 4.0\times 10^{2}\,\mathrm{rad/s}
\]

对应的衰减时间：
\[
\tau = \frac{1}{|\omega_i|} \approx \frac{1}{4.0\times 10^{2}} \approx 2.5\times 10^{-3}\,\mathrm{s} = 2.5\,\mathrm{ms}
\]

\textbf{分析}：$|\omega_i|/\omega_{pe}\approx7\times10^{-9}$，远小于 $\omega_r$，弱阻尼假设自洽。波在毫秒尺度衰减，说明\textbf{长波} Langmuir 波在热等离子体中几乎无阻尼。这一结果对波诊断有直接意义：想用探针或散射测量 Langmuir 波，必须选 $k\lambda_D\ll1$ 的长波，否则波自身还没传播开就被阻尼掉了。

\textbf{越界会怎样（提高题）：} 把 $k$ 换成 $5\times10^{4}\,\mathrm{m^{-1}}$（$k\lambda_D=0.745$），代入同一公式会得到 $|\omega_i|\approx7.7\times10^{9}\,\mathrm{rad/s}$、$\tau\approx0.13\,\mathrm{ns}$。\textbf{但这个数不可信}——$k\lambda_D=0.745$ 已不在渐近区（要求 $\ll1$），$\exp[-1/(2k^2\lambda_D^2)]$ 的展开前提（相速度落在分布尾部）已被破坏，正确做法是数值求解复色散函数 $Z(\zeta)$ 的根。请自行验证：当 $k\lambda_D\to1$ 时，用渐近式与用 $Z$ 函数求得的阻尼率分歧迅速增大；只有当 $k\lambda_D\lesssim0.3$ 时两者才吻合。这正是本书强调"渐近式必须标出适用范围"的原因。
\end{example}

\begin{example}{共振粒子的能量交换——Landau阻尼的微观机制}{resonant-energy}
设一维电子麦克斯韦分布的热速度\emph{统一}定义为 $v_{te} = \sqrt{k_{\!B}T_e/m_e}$（与本书 $\lambda_{De}=v_{te}/\omega_{pe}$ 的定义一致），Langmuir波的相速度为 $v_\phi = 1.5\,v_{te}$。取共振区间宽度为 $\Delta v = 0.1\,v_{te}$，估算在 $v_\phi - \Delta v < v < v_\phi + \Delta v$ 范围内，慢粒子（$v < v_\phi$）与快粒子（$v > v_\phi$）的数量差异，并解释为什么这导致净能量从波流向粒子。

\textbf{约定提示：} 全题只用 $v_{te}=\sqrt{k_{\!B}T_e/m_e}$ 一种定义，$v_\phi=1.5\,v_{te}$ 与 $\Delta v=0.1\,v_{te}$ 都相对它取。请勿中途换成 $v_t=\sqrt{2k_{\!B}T_e/m_e}$——那会使"1.5"和"0.1"实际变成一个不同的物理速度。若要换定义，两个无量纲比值必须同时按 $\sqrt2$ 重标。
\tcblower
\textbf{解答：}

一维麦克斯韦分布（$v_{te}=\sqrt{k_{\!B}T_e/m_e}$）：
\[
f(v) = \frac{n_0}{\sqrt{2\pi}\,v_{te}}\,\exp\!\left(-\frac{v^2}{2v_{te}^2}\right)
\]

\textbf{步骤1：相速度处的分布函数与斜率。} 令 $x\equiv v/v_{te}$，则 $f\propto\ee^{-x^2/2}$，其对数斜率为
\[
\frac{f'(v)}{f(v)} = -\frac{v}{v_{te}^2}
\quad\Rightarrow\quad
f'(v_\phi) = -\frac{v_\phi}{v_{te}^2}f(v_\phi) = -\frac{1.5}{v_{te}}f(v_\phi)\cdot\frac{1}{v_{te}}
\]
（注意 $f'$ 的量纲是 $f$ 除以速度：$[f']=[n_0]/[\mathrm{m/s}]^2$，所以斜率里\emph{恰好一个} $1/v_{te}$，不是两个也不是零。）

\textbf{步骤2：慢、快区间的粒子数。} 对窄区间 $\Delta v\ll v_{te}$ 用线性近似 $f(v)\approx f(v_\phi)+f'(v_\phi)(v-v_\phi)$：
\[
N_{<} = \int_{v_\phi-\Delta v}^{v_\phi}\!\! f\,\diff v \approx f(v_\phi)\Delta v - \tfrac12 f'(v_\phi)\Delta v^2,
\qquad
N_{>} = \int_{v_\phi}^{v_\phi+\Delta v}\!\! f\,\diff v \approx f(v_\phi)\Delta v + \tfrac12 f'(v_\phi)\Delta v^2
\]

\textbf{步骤3：数量差异与相对差异。}
\[
N_{<}-N_{>} = -f'(v_\phi)\Delta v^2 = \frac{v_\phi}{v_{te}^2}f(v_\phi)\Delta v^2
\]
\[
\frac{N_{<}-N_{>}}{N_{<}+N_{>}} \approx \frac{-f'(v_\phi)\Delta v^2}{2f(v_\phi)\Delta v}
= \frac{v_\phi\,\Delta v}{2\,v_{te}^2}
= \frac{(1.5v_{te})(0.1v_{te})}{2v_{te}^2}
= \frac{0.15}{2} = \boxed{7.5\%}
\]

\textbf{关键：} 比例里 $v_{te}^2$ 精确抵消，结果\emph{只}依赖两个无量纲比值 $v_\phi/v_{te}=1.5$ 和 $\Delta v/v_{te}=0.1$，所以\textbf{与选哪种热速度定义无关}——前提是两个比值相对\emph{同一种}定义给出。数值核验：对 $\ee^{-v^2/2}$ 直接在 $[1.4,1.6]$ 上积分，得 $(N_<-N_>)/(N_<+N_>)=0.0748\approx7.5\%$，与线性近似吻合。

\textbf{物理解释：}
\begin{itemize}[leftmargin=*,itemsep=0pt]
\item 慢粒子（$v < v_\phi$）在波参考系中向左运动（$v-v_\phi<0$），被波的电场加速，从波获得动能。由于 $\partial f_0/\partial v<0$，慢粒子比快粒子多。
\item 快粒子（$v > v_\phi$）在波参考系中向右运动，被波的电场减速，把动能交给波。
\item 因为慢粒子更多，净效应是粒子从波获得能量——波被阻尼。这就是 Landau 阻尼的微观图像：不是碰撞，而是速度空间分布的斜率决定能量流向。
\end{itemize}

若分布有正斜率（$\partial f_0/\partial v>0$，如 bump-on-tail），则快粒子更多，净能量从粒子流向波——波增长，即不稳定性。
\end{example}

\begin{example}{bump-on-tail不稳定性条件判断}{bump-on-tail-instability}
某等离子体中，背景电子密度为 $n_b = 10^{18}\,\mathrm{m}^{-3}$，温度为 $T_b = 4\,\mathrm{eV}$。另有一束高能电子，密度 $n_b = 10^{16}\,\mathrm{m}^{-3}$（为背景密度的 $1\%$），漂移速度 $u = 5\times 10^{7}\,\mathrm{m/s}$，束流温度 $T_s = 0.5\,\mathrm{eV}$。判断该体系是否满足bump-on-tail不稳定性条件。
\tcblower
\textbf{解答：}

\textbf{步骤1：计算背景电子热速度。}
\[
v_{tb} = \sqrt{\frac{k_{\!B}T_b}{m_e}}
= \sqrt{\frac{6.408\times 10^{-19}}{9.109\times 10^{-31}}}
\approx 8.39\times 10^{5}\,\mathrm{m/s} \approx 8.4\times 10^{5}\,\mathrm{m/s}
\]

\textbf{步骤2：计算束流热速度。}
\[
v_{ts} = \sqrt{\frac{k_{\!B}T_s}{m_e}}
= \sqrt{\frac{0.5\times 1.602\times 10^{-19}}{9.109\times 10^{-31}}}
\approx 2.97\times 10^{5}\,\mathrm{m/s}
\]

\textbf{步骤3：判断不稳定性条件。}

bump-on-tail不稳定性的核心是：在相速度 $v_\phi = \omega/k$ 处，分布函数有正斜率（$\partial f/\partial v > 0$）。这要求束流的漂移速度 $u$ 足够大，使得束流峰与背景麦克斯韦分布的尾部之间形成凸起。

定性判据：束流速度 $u$ 应远大于背景热速度，但束流自身热速度 $v_{ts}$ 应足够小，使束流峰不"摊平"分布。

定量上，常用条件是：
\[
\frac{u}{v_{tb}} \gtrsim 3\quad\text{且}\quad
\frac{n_s}{n_b}\left(\frac{u}{v_{tb}}\right)^2 \gtrsim 1
\]

检查：
\[
\frac{u}{v_{tb}} = \frac{5\times 10^{7}}{8.4\times 10^{5}} \approx 59.5 \gg 3 \quad\checkmark
\]
\[
\frac{n_s}{n_b}\left(\frac{u}{v_{tb}}\right)^2
= 0.01\times (59.5)^2 \approx 0.01\times 3540 \approx 35.4 \gg 1 \quad\checkmark
\]

\textbf{结论}：两个条件均强烈满足。该体系必然存在 bump-on-tail 不稳定性。束流电子在速度空间中形成明显的"凸起"，在 $v \approx u$ 附近 $\partial f/\partial v > 0$，Langmuir 波将从束流获得能量而增长。

\textbf{补充}：不稳定性的最大增长率约为
\[
\gamma_{\max} \sim \frac{n_s}{2n_b}\omega_{pe}
\]
其中 $\omega_{pe} \approx 5.64\times 10^{10}\,\mathrm{rad/s}$，故 $\gamma_{\max} \sim 0.005\times 5.64\times 10^{10} \approx 2.8\times 10^{8}\,\mathrm{rad/s}$，对应增长时间约 $3.6\,\mathrm{ns}$。
\end{example}

\begin{figure}[htbp]
\centering
\begin{tikzpicture}[scale=1.1,font=\small]
% 坐标轴
\draw[->,thick] (-0.5,0) -- (6.5,0) node[right] {$v$};
\draw[->,thick] (0,-0.3) -- (0,4.5) node[above] {$f(v)$};

% 麦克斯韦分布曲线
\draw[domain=-0.5:5.5,samples=100,thick,blue] plot({\x+0.5},{3.8*exp(-(\x-2)^2/1.5)});

% 相速度标注
\draw[dashed,red] (3.5,0) -- (3.5,3.2);
\node[red,below] at (3.5,-0.2) {$v_\phi$};

% 共振区间
\draw[decorate,decoration={brace,amplitude=5pt},thick,red!70] (3.1,3.4) -- (3.9,3.4);
\node[red!70,above] at (3.5,3.6) {\footnotesize 共振区};

% 慢粒子区标注
\fill[blue!20] (2.5,0) rectangle (3.5,2.8);
\node[blue!60!black,align=center] at (2.3,1.5) {\footnotesize 慢粒子多\\\footnotesize $v<v_\phi$};

% 快粒子区标注
\fill[orange!20] (3.5,0) rectangle (4.5,1.5);
\node[orange!60!black,align=center] at (4.7,0.8) {\footnotesize 快粒子少\\\footnotesize $v>v_\phi$};

% 斜率示意
\draw[->,thick,green!50!black] (3.5,2.2) -- (2.8,2.6);
\node[green!50!black,right] at (3.6,2.3) {\footnotesize $\partial f/\partial v < 0$};

% 能量流向
\draw[->,very thick,red!70] (3.5,-1) -- (2.5,-1);
\node[red!70,below] at (3,-1.3) {\footnotesize 净能量流向粒子};
\end{tikzpicture}
\caption{Landau阻尼的共振粒子示意图。在相速度 $v_\phi$ 附近，速度略小于 $v_\phi$ 的粒子比速度略大于 $v_\phi$ 的粒子多（$\partial f_0/\partial v < 0$），导致净能量从波流向粒子，波被阻尼。}
\end{figure}

\begin{figure}[htbp]
\centering
\begin{tikzpicture}[scale=1.1,font=\small]
% 坐标轴
\draw[->,thick] (-0.5,0) -- (8,0) node[right] {$v$};
\draw[->,thick] (0,-0.3) -- (0,4.5) node[above] {$f(v)$};

% 背景麦克斯韦分布（大宽峰）
\draw[domain=-0.5:5,samples=100,thick,blue] plot({\x+0.5},{3.5*exp(-(\x-1.5)^2/1.8)});

% 束流峰（小窄峰）
\draw[domain=4.5:7,samples=50,thick,red] plot({\x},{1.8*exp(-(\x-5.8)^2/0.3)});

% 合成分布（示意）
\draw[domain=-0.5:5,samples=100,thick,dashed,gray] plot({\x+0.5},{3.5*exp(-(\x-1.5)^2/1.8)+0.1});
\draw[domain=4.5:7,samples=50,thick,dashed,gray] plot({\x},{1.8*exp(-(\x-5.8)^2/0.3)+0.1});

% 标注
\node[blue,align=center] at (1.5,3.8) {\footnotesize 背景等离子体};
\node[blue,align=center] at (1.5,3.4) {\footnotesize $n_b$, $T_b$};

\node[red,align=center] at (6.2,2.3) {\footnotesize 电子束流};
\node[red,align=center] at (6.2,1.9) {\footnotesize $n_s$, $u$, $T_s$};

% 凸起区域标注
\draw[->,thick,green!50!black] (4.5,1.5) -- (5.5,2.2);
\node[green!50!black,align=center] at (4.2,1.2) {\footnotesize bump（凸起）};
\node[green!50!black,align=center] at (4.2,0.8) {\footnotesize $\partial f/\partial v > 0$};

% 不稳定性区域
\fill[pattern=north east lines,pattern color=green!50] (4.2,0) rectangle (5.5,0.5);
\node[green!50!black] at (4.85,-0.4) {\footnotesize 不稳定区};
\end{tikzpicture}
\caption{bump-on-tail速度分布函数示意图。蓝色宽峰为背景热电子，红色窄峰为漂移电子束流。在两峰之间的"尾部"区域形成凸起（bump），对应 $\partial f/\partial v > 0$，驱动Langmuir波不稳定性增长。}
\end{figure}

% ch7 新例题：电子等离子体波的朗道阻尼数值练习
\begin{example}{朗道阻尼随波长的变化}{landau-wavelength-dependence}
氩等离子体（$n_e=10^{18}\,\mathrm{m^{-3}}$，$T_e=4\,\mathrm{eV}$）中，比较波数 $k = 10^4\,\mathrm{m^{-1}}$ 和 $k = 10^5\,\mathrm{m^{-1}}$ 的两列朗缪尔波的朗道阻尼率，并讨论波长变短时阻尼如何变化。
\tcblower
\textbf{解答：}

由例题 7.2 的方法：$\lambda_D = 1.49\times10^{-5}\,\mathrm{m}$（$14.9\,\mu\mathrm{m}$），$\omega_{pe} = 5.64\times10^{10}\,\mathrm{rad/s}$，朗道阻尼率公式
\[
\omega_i = -\sqrt{\frac{\pi}{8}}\,\frac{\omega_{pe}}{(k\lambda_D)^3}\exp\!\left(-\frac{1}{2k^2\lambda_D^2}-\frac{3}{2}\right)
\]

\textbf{情形 1：$k = 10^4\,\mathrm{m^{-1}}$。} $k\lambda_D = 0.149$：
\[
\frac{1}{2(k\lambda_D)^2} = \frac{1}{2\times0.0222} = 22.5,\qquad \text{指数} = -22.5-1.5 = -24
\]
\[
|\omega_i| = 0.627\times\frac{5.64\times10^{10}}{(0.149)^3}\times\ee^{-24} \approx 3.9\times10^{2}\,\mathrm{rad/s}
\]
衰减时间 $\tau = 1/|\omega_i| \approx 2.5\,\mathrm{ms}$——\textbf{极弱阻尼}，波可以传播很远。

\textbf{情形 2：$k = 10^5\,\mathrm{m^{-1}}$。} $k\lambda_D = 1.49$。\textbf{这一情形不能直接套用上面的渐近式}：$k\lambda_D=1.49>1$，公式的展开前提已失效。若强行代入，会得到 $|\omega_i|\approx1.9\times10^{9}\,\mathrm{rad/s}$（注意与例题7.2原稿的 $7.8\times10^{9}$ 不同——那个数是照抄 $k\lambda_D=0.745$ 的情形，与前式代数不符）。正确结论应表述为：\textbf{此时必须数值求解含 $Z(\zeta)$ 的复色散关系}，渐近式给出的量级仅供定性参考。

\textbf{规律：} 三个波数对应 $k\lambda_D = 0.149,0.745,1.49$。在\emph{渐近式仍可信}的区间内（约 $k\lambda_D\lesssim0.3$），$k\lambda_D$ 从 $0.149$ 略增，指数因子 $\ee^{-24}$ 已经极小，阻尼率随 $k\lambda_D$ 增大而急剧上升——指数中 $1/(2k^2\lambda_D^2)$ 对 $k\lambda_D$ 极为敏感。朗道阻尼的物理根源是：长波（$k\lambda_D\ll1$）的相速度远大于电子热速度，共振粒子数按 $\exp[-v_\phi^2/(2v_t^2)]$ 指数稀少；波长缩短使相速度落向电子热分布的主体，共振粒子迅速增多，阻尼随之剧增。

\textbf{本例题的正确用法：} 情形 1 用于展示有效的长波估算；情形 2 用于说明\textbf{何时不能套公式}。要在 $k\lambda_D\sim1$ 处得到可信的阻尼率，需解复色散函数（Fitzpatrick《Plasma Physics》第 5 章有标准处理）。这正是朗缪尔探针测量电子分布函数（Druyvesteyn法，第12章）与波诊断选波长时必须考虑的基本限制。
\end{example}

\begin{formulabox}{第7章核心公式速查}{ch7-formulas}
\begin{itemize}[leftmargin=*,itemsep=0pt]
\item Vlasov方程：$\displaystyle\frac{\partial f}{\partial t} + \vect{v}\cdot\grad f + \frac{q}{m}(\vect{E} + \vect{v}\times\vect{B})\cdot\grad_v f = 0$
\item 线性化色散关系（静电波）：$\displaystyle 1 = \frac{\omega_{pe}^2}{n_0k^2}\int_{-\infty}^{\infty} \frac{\partial f_0/\partial v}{v - \omega/k}\,\diff v$，其中 $\int f_0\diff v=n_0$（数密度分布约定；若用归一化 $\hat f_0=f_0/n_0$ 则去掉 $1/n_0$）
\item Landau阻尼率（麦克斯韦分布，\textbf{长波弱阻尼渐近式}，须 $k\lambda_D\ll1$）：$\displaystyle\omega_i = -\sqrt{\frac{\pi}{8}}\,\frac{\omega_{pe}}{(k\lambda_D)^3}\exp\!\left(-\frac{1}{2k^2\lambda_D^2} - \frac{3}{2}\right)$
\item 波相速度：$\displaystyle v_\phi = \frac{\omega}{k}$
\item 共振条件：$\displaystyle v_{\text{res}} = v_\phi = \frac{\omega}{k}$
\item bump-on-tail增长率（估算）：$\displaystyle\gamma_{\max} \sim \frac{n_s}{2n_b}\omega_{pe}$
\item 德拜长度：$\displaystyle\lambda_D = \sqrt{\frac{\eps_0 k_{\!B}T_e}{n_e e^2}}$
\item 电子等离子体频率：$\displaystyle\omega_{pe} = \sqrt{\frac{n_e e^2}{\eps_0 m_e}}$
\end{itemize}
\end{formulabox}

\begin{checklistbox}{第7章自学检查清单}{ch7-checklist}
\begin{itemize}[leftmargin=*,label=$\square$]
\item 我能写出Vlasov方程，并解释其物理意义（相空间不可压缩流、刘维尔定理）
\item 我能对Vlasov方程进行线性化，得到色散关系中的极点项
\item 我理解Landau围道的必要性：为什么必须从极点下方绕过（因果律）
\item 我能从速度分布的斜率判断Landau阻尼（$\partial f_0/\partial v < 0$）或不稳定性（$\partial f_0/\partial v > 0$）
\item 我能估算给定 $k\lambda_D$ 下的Landau阻尼率
\item 我理解bump-on-tail不稳定性的物理机制：束流在速度空间形成凸起，共振粒子向波传递能量
\item 我知道Landau阻尼是线性效应，而粒子 trapping 是非线性效应
\end{itemize}
\end{checklistbox}

\endinput
